
\documentclass{article}
\usepackage{colortbl}
\usepackage{makecell}
\usepackage{multirow}
\usepackage{supertabular}

\begin{document}

\newcounter{utterance}

\centering \large Interaction Transcript for game `eqbench', experiment `eq\_bench', episode 143 with Qwen3.6{-}27B.
\vspace{24pt}

{ \footnotesize  \setcounter{utterance}{1}
\setlength{\tabcolsep}{0pt}
\begin{supertabular}{c@{$\;$}|p{.15\linewidth}@{}p{.15\linewidth}p{.15\linewidth}p{.15\linewidth}p{.15\linewidth}p{.15\linewidth}}
    \# & \multicolumn{2}{c}{Player} && \multicolumn{2}{c}{Game Master} \\
    \hline

    \theutterance \stepcounter{utterance}  
    & & & \multicolumn{4}{p{0.6\linewidth}}{
        \cellcolor[rgb]{0.9,0.9,0.9}{
            \makecell[{{p{\linewidth}}}]{
                \texttt{\tiny{[P1$\langle$GM]}}
                \\ 
\texttt{Your task is to predict the likely emotional responses of a character in this dialogue:} \\
\\ 
\texttt{Amelia: I find it amusing, Xarlon, that you claim your civilization is superior, yet you are here, on Earth, seeking dialogue with us.} \\
\texttt{Xarlon: Superiority is not a matter of isolation, Dr. Singh. It is about understanding, learning, and evolving.} \\
\texttt{Amelia: Evolving? Is that what you call it? Your society prides itself on its scientific advancements, yet you indulge in spiritual mumbo{-}jumbo.} \\
\texttt{Xarlon: You belittle what you do not understand, Dr. Singh. Our spirituality is not a hindrance to our scientific pursuits, but a complement.} \\
\texttt{Amelia: A complement? Spirituality is the antithesis of science. It's unquantifiable, unprovable. It's a crutch for the weak{-}minded.} \\
\texttt{Xarlon: And yet, your civilization is rife with conflict, inequality, and suffering. Could it be that your rejection of the intangible has led you astray?} \\
\texttt{Amelia: We are not perfect, Xarlon, but we don't pretend to be. We acknowledge our flaws, our mistakes. We learn. We grow.} \\
\texttt{Xarlon: And yet, your progress is slow. Your planet is dying, your people divided. Is this the result of your 'learning'?} \\
\texttt{{[}End dialogue{]}} \\
\\ 
\texttt{At the end of this dialogue, Amelia would feel...} \\
\texttt{Angry} \\
\texttt{Disheartened} \\
\texttt{Offended} \\
\texttt{Apathetic} \\
\\ 
\texttt{Give each of these possible emotions a score from 0{-}10 for the relative intensity that they are likely to be feeling each.} \\
\\ 
\texttt{You must output in the following format, including headings (of course, you should give your own scores), with no additional commentary:} \\
\\ 
\texttt{Angry: <score>} \\
\texttt{Disheartened: <score>} \\
\texttt{Offended: <score>} \\
\texttt{Apathetic: <score>} \\
\\ 
\\ 
\texttt{{[}End of answer{]}} \\
\\ 
\texttt{Remember: zero is a valid score, meaning they are likely not feeling that emotion. You must score at least one emotion > 0.} \\
\\ 
\texttt{Your answer:} \\
            }
        }
    }
     \\ \\

    \theutterance \stepcounter{utterance}  
    & \multicolumn{4}{p{0.6\linewidth}}{
        \cellcolor[rgb]{0.561,0.737,0.561}{
            \makecell[{{p{\linewidth}}}]{
                \texttt{\tiny{[P1$\rangle$GM]}}
                \texttt{Angry: 8} \\
\texttt{Disheartened: 4} \\
\texttt{Offended: 7} \\
\texttt{Apathetic: 0} \\
            }
        }
    }
    & & \\ \\

    \theutterance \stepcounter{utterance}  
    & & & \multicolumn{2}{p{0.3\linewidth}}{
        \cellcolor[rgb]{0.9,0.9,0.9}{
            \makecell[{{p{\linewidth}}}]{
                \texttt{\tiny{[GM$|$GM]}}
                \texttt{Angry: 8} \\
\texttt{Disheartened: 4} \\
\texttt{Offended: 7} \\
\texttt{Apathetic: 0} \\
            }
        }
    }
    & & \\ \\

    \theutterance \stepcounter{utterance}  
    & & & \multicolumn{2}{p{0.3\linewidth}}{
        \cellcolor[rgb]{0.9,0.9,0.9}{
            \makecell[{{p{\linewidth}}}]{
                \texttt{\tiny{[GM$|$GM]}}
                \texttt{Angry: 7} \\
\texttt{Disheartened: 0} \\
\texttt{Offended: 6} \\
\texttt{Apathetic: 0} \\
            }
        }
    }
    & & \\ \\

    \theutterance \stepcounter{utterance}  
    & & & \multicolumn{2}{p{0.3\linewidth}}{
        \cellcolor[rgb]{0.9,0.9,0.9}{
            \makecell[{{p{\linewidth}}}]{
                \texttt{\tiny{[GM$|$GM]}}
                \texttt{game\_result = LOSE} \\
            }
        }
    }
    & & \\ \\

\end{supertabular}
}

\end{document}
